\documentclass[12pt]{article}
\usepackage[T1]{fontenc}
\usepackage[utf8]{inputenc}
\usepackage{lmodern}
\usepackage{textcomp}
\usepackage{enumitem}
\usepackage{geometry}
\geometry{margin=1in}
\begin{document}
\begin{center}
Answer Key - Business Administration - Undergraduate Level Assessment 1 \
\small (Answers are ordered exactly as in the question paper.)
\end{center}

\section*{Multiple Choice}
\begin{enumerate}[label=\arabic*.]
\item \textbf{Answer:} A
\\ \textit{Options:} A) They offer flexibility of talents and skills.; B) They bring greater credibility to the work.; C) They will provide sound ROI on the project or program.; D) They can provide more attention and detail to a project that in-house staff.
\\ \textit{Note:} Organizations retain outside public relations counsel because these professionals provide a diverse range of specialized skills and expertise that can be tailored to meet the specific and evolving needs of the organization, thereby offering greater flexibility compared to in-house teams.
\item \textbf{Answer:} B
\\ \textit{Options:} A) Media release; B) Media tour; C) Press room; D) Promotional days/weeks
\\ \textit{Note:} A media tour is a public relations tactic that specifically involves arranging visits for journalists to relevant locations to provide them with firsthand experiences and insights for their stories.
\item \textbf{Answer:} A
\\ \textit{Options:} A) Direct.; B) Indirect.; C) Hybrid.; D) None of the above.
\\ \textit{Note:} The correct answer is "Direct" because in a direct channel structure, the producer sells their products straight to the final customer without intermediaries, ensuring a direct relationship and streamlined distribution.
\item \textbf{Answer:} B
\\ \textit{Options:} A) Encoding.; B) Decoding.; C) Transfer.; D) Noise.
\\ \textit{Note:} Decoding is the process whereby receivers interpret and make sense of the incoming message, transforming it from its encoded form back into comprehensible information.
\item \textbf{Answer:} D
\\ \textit{Options:} A) Worker welfare; B) Health and safety; C) Interpersonal relationships; D) Productivity
\\ \textit{Note:} The 'rational goal' model of management emphasizes achieving specific organizational objectives through efficient processes and high productivity levels, making productivity the most commonly associated characteristic.
\end{enumerate}

\section*{True / False}
\begin{enumerate}[label=\arabic*.]
\setcounter{enumi}{5}
\item \textbf{Answer:} True
\\ \textit{Options:} A) True; B) False
\\ \textit{Note:} The statement is true because marketing fundamentally revolves around understanding and meeting customer needs and preferences, thereby prioritizing customer value and satisfaction as core elements of successful business strategies.
\item \textbf{Answer:} True
\\ \textit{Options:} A) True; B) False
\\ \textit{Note:} Examining public records is considered a key component of secondary research because it allows researchers to gather existing data and insights from previously collected information, which can inform their analysis and decision-making processes in business administration.
\item \textbf{Answer:} False
\\ \textit{Options:} A) True; B) False
\\ \textit{Note:} A span of control of 5 is not universally considered the most effective because the optimal span of control can vary significantly based on factors such as the nature of the tasks, the complexity of the organization, and the capabilities of the manager and employees involved.
\item \textbf{Answer:} False
\\ \textit{Options:} A) True; B) False
\\ \textit{Note:} Global production networks increasingly emphasize sustainability and waste reduction as integral to their profit strategies, recognizing that efficient waste recapture can lead to cost savings and enhanced brand reputation.
\item \textbf{Answer:} True
\\ \textit{Options:} A) True; B) False
\\ \textit{Note:} Henry (2006) argues that a well-defined strategy enables organizations to effectively allocate resources, make informed decisions, and differentiate themselves from competitors, thereby securing a competitive advantage.
\end{enumerate}

\section*{Long Form Response}
\begin{enumerate}[label=\arabic*.]
\setcounter{enumi}{10}
\item \textbf{Answer:} Buycotts play a significant role in promoting favorable behavior among companies by encouraging consumers to support businesses that align with their values, thus incentivizing those companies to adopt socially responsible practices. This consumer- driven approach can lead to increased sales as companies respond to the positive reinforcement of their initiatives. Digital technology enhances the effectiveness of buycotts by facilitating communication and mobilization through social media, allowing consumers to quickly organize and share information about their preferred brands. Additionally, online platforms enable tracking of sales and consumer engagement, making it easier for companies to gauge the impact of their socially responsible actions. As a result, the intersection of buycotts and digital technology not only drives increased sales but also fosters a corporate culture that prioritizes ethical behavior.
\\ \textit{Note:} correct because it illustrates how buycotts leverage consumer values to drive corporate responsibility and increase sales, while emphasizing the role of digital technology in amplifying awareness and mobilization for these campaigns.
\item \textbf{Answer:} Muckraking journalism during the Seedbed Era (1900-1917) played a crucial role in exposing social injustices and corporate malpractices, significantly influencing public perception and trust. Journalists like Upton Sinclair and Ida Tarbell highlighted issues such as labor exploitation and corrupt business practices, leading to heightened public awareness and demand for reform. This scrutiny compelled businesses to adopt more transparent and ethical communication strategies, laying the groundwork for modern public relations. As companies recognized the importance of managing their image in response to muckraking investigations, they began to implement proactive public relations practices to foster goodwill and mitigate negative publicity. Ultimately, muckraking journalism not only shaped public discourse but also transformed how businesses approached stakeholder communication and reputation management.
\\ \textit{Note:} correct because it emphasizes how muckraking journalism revealed societal and corporate issues that fostered public demand for accountability, prompting businesses to embrace transparency and ethical communication, which are fundamental principles in the evolution of public relations.
\end{enumerate}

\vfill
\noindent\textit{End of Answer Key}
\end{document}